\begin{lemma}[the primitive circle marking]\label{push:src:C.3.9}
Use the conductor generators \(a,b\) of \ref{push:src:C.3.7}, with
\(E_1\) represented by the edge of direction \(e_1-e_2\),
\(E_2\) by the edge of direction \(e_1\), and the closing path
\(c\) on the fixed edge of direction \(e_2\). Under the
natural period marking of the cusp orbit fibre,
\[
 a\longmapsto\widehat u,\qquad
 b\longmapsto\widehat\gamma+\widehat u.
\]
In particular the circle obtained by gluing the two fixed boundary
spheres of the normalization has oriented generator
\(ab^{-1}=-\widehat\gamma\). It is primitive and has no
additional \(\widehat u\) component.
\end{lemma}
\begin{proof}
The cusp period action has fan translation
\(\lambda\mapsto B_0\lambda\), with
\[
 B_0=\begin{pmatrix}0&1\\-1&0\end{pmatrix},
 \quad \lambda=(\gamma,u)^t.
\]
This is the actual period matrix
\(Z=sB_0+C(q)\) of the cusp construction, not an unmarked
torus identification. The unit factors \(\exp(2\pi i C(q)\lambda)\)
fix the toric strata and hence do not change their triangle labels.
Take the cyclic fan directions
\[
 (D_1,\ldots,D_6)=
 (e_1-e_2,e_1,e_2,e_2-e_1,-e_1,-e_2).
\]
For \(L_v=(v,v+e_1,v+e_2)\) and
\(U_v=(v+e_1,v+e_2,v+e_1+e_2)\), the vertices of the
normalizing component have the following actual lifts:
\[
 \begin{array}{c|cccccc}
 &p_1&p_2&p_3&p_4&p_5&p_6\\\hline
 \text{triangle}&U_{-e_2}&L_0&U_{-e_1}&L_{-e_1}
                 &U_{-e_1-e_2}&L_{-e_2}.
 \end{array}
\]
For example \(p_1\) is the triangle
\(\{0,e_1-e_2,e_1\}=U_{-e_2}\); the other entries follow
by the same explicit substitution.

Choose \(A=L_0\) and the lift of \(c\) from
\(p_2=L_0\) to \(p_3=U_{-e_1}\).
The edge \(D_1\), oriented from \(A\) to \(B\), goes
from \(L_{-e_2}\) to \(U_{-e_2}\). Translate it by
\(+e_2\) to start at the selected lift \(L_0\).
It ends at \(U_0\), which is the chosen \(B\)-lift
translated by \(+e_1\). Thus \(a=E_1c^{-1}\) has
fan displacement \(+e_1=B_0\widehat u\).
The edge \(D_2\) oriented from \(A\) to \(B\)
starts already at \(L_0\) and ends at \(U_{-e_2}\).
Its displacement relative to the chosen \(B\)-lift is
\(e_1-e_2=B_0(\widehat\gamma+\widehat u)\).
This proves the two labels.

The real normalization circle is the path along \(D_1\)
and back along \(D_2\), closed by the gluing of the fixed
boundary spheres. Its image is \(E_1E_2^{-1}=ab^{-1}\),
so its label is \(-\widehat\gamma\). Equivalently, the
fixed-edge gluing has fan translation
\(B_0(1,0)^t=-e_2\), corresponding to the period
\(\widehat\gamma\); traversing its fundamental
domain in the stated direction has the opposite sign.
The two computed labels have determinant \(-1\), and hence
give the entire period lattice. A loop lift is legitimate even
at a vertex: a nonzero fan translation cannot preserve the finite
set of component labels incident to that vertex, so the period
action on the toric cover, also after the circle quotient, is free.
The finite triangle and displacement calculation is checked in
\path{code/d_push_cusp_marking.py}.
\end{proof}

